\documentclass[12pt]{article}

\usepackage{amsmath}
\usepackage{amssymb}
\usepackage{graphicx}
\graphicspath{{figures/}{figures/ic_search/}}
\usepackage{hyperref}
\usepackage[margin=1in]{geometry}
\usepackage{xcolor}
\usepackage{booktabs}
\usepackage{natbib}
\bibliographystyle{aasjournal}

\newcommand{\bk}{\mathbf{k}}
\newcommand{\bx}{\mathbf{x}}
\newcommand{\avg}[1]{\left\langle #1 \right\rangle}
\newcommand{\Mpch}{\ \mathrm{Mpc}/h}

\title{Selecting an Initial-Conditions Seed on the Power Spectrum\\
       of a Subvolume}
\author{}
\date{\today}

\begin{document}
\maketitle

\begin{abstract}
A subvolume of a simulation box does not measure the linear power spectrum. Its
estimator is unbiased for the theory convolved with the subvolume's own window,
which for the geometry used here lies $34\%$ below the theory at the fundamental
mode. Choosing an initial-conditions seed because its subvolume spectrum matches
the unconvolved theory therefore selects realisations whose large-scale power is
anomalously high, by construction rather than by accident. We measure what this
does using $83\,542$ independent realisations at $N = 128^3$,
$L = 700\Mpch$. Selection displaces the subvolume's large-scale power by
$1.26\sigma$ and every other property by its correlation with that power times
$1.26$: the tidal shear by $0.62\sigma$, the cosmic-web composition by
$0.35\sigma$, the mean density, whose correlation is $0.001$, not at all. The
displacement saturates, because the realisation scatter bounds how far the power
can be pushed, so a search over a million seeds is no worse than one over a
thousand. Two controls confirm the cause: a randomly permuted criterion produces
$0.000 \pm 0.045$ in every property, and selecting against the convolved theory,
which asks for nothing atypical, produces $0.016$. Fitting an amplitude with the
window in the model returns $A = 0.9995 \pm 0.0765$ over all subvolumes with no
search at all, while the selected subvolumes have $A = 1.065$. The result is not
specific to power spectra: any rule for choosing which region to resimulate
conditions the field, and the correlation between the rule and the quantity of
interest predicts the bias in advance.
\end{abstract}

\section{The proposal}
\label{sec:proposal}

Generate initial conditions with many random seeds. In each, measure the power
spectrum inside the subvolume that will be resimulated at higher resolution.
Adopt the seed whose subvolume spectrum best matches linear theory, on the
grounds that such a region is a faithful piece of the universe to zoom into.

The intent is reasonable, and the failure is not obvious from the outside. This
note sets out why the criterion cannot do what it is meant to do, measures how
large the resulting distortion is, and shows that the distortion is bounded and
predictable.

Throughout, the subvolume is a pencil: one eighth of the box across in two
directions and its full length in the third, so
$87.5 \times 87.5 \times 700\ (\Mpch)^3$ in a $700\Mpch$ box. Fields are
Lagrangian densities $\delta(q)$ at $z = 200$ from monofonIC with second-order
Lagrangian perturbation theory, on a $128^3$ grid.

\section{A subvolume does not measure the theory}
\label{sec:window}

Restricting a field to a region multiplies it by a mask. Multiplication in
configuration space is convolution in Fourier space, so the power spectrum
measured in the region is not the theory but the theory smeared with the
region's window,
\begin{equation}
  \label{eq:conv}
  \avg{\hat P(\bk)} \;=\; \big[P_{\rm th} \ast |\widetilde W|^2\big](\bk)
  \;\neq\; P_{\rm th}(\bk).
\end{equation}
The smearing moves power from the largest scales, where the spectrum is
steepest, into neighbouring wavenumbers. The effect is deterministic. It is
identical in every realisation, it does not average away with more
realisations, and it is not noise. For this geometry it is a $34\%$ deficit at
the fundamental mode.

Equation~\eqref{eq:conv} is a statement about an average, so it can be tested
directly. Averaging $2\,005\,008$ subvolume measurements against the convolved
curve gives $0.9996 \pm 0.0001$, and against the full box the unconvolved theory
gives $0.99994 \pm 0.00002$. The estimator is unbiased for the convolved theory
and biased against the raw theory, exactly as the equation says.

The measurement is therefore not failing to reach the theory. It is succeeding
at measuring something else: the theory as seen through a window of that size
and shape.

\section{What the search selects}
\label{sec:selects}

If the mean measurement lies $34\%$ below a curve, the only realisations that
reach that curve are those with anomalously high power. A search for a matching
seed does not find a region that agrees with the theory; it finds a region
unusual in precisely the way needed to cancel a geometric deficit.

The cleanest way to see this is to stop searching and fit instead. For each
subvolume, fit the amplitude $A$ in $\hat P(k) = A\,P_{\rm model}(k)$ by
generalised least squares using the covariance measured from the ensemble, once
with the window in the model and once without. Table~\ref{tab:amplitude} gives
the result and Figure~\ref{fig:amplitude} shows the distributions.

\begin{table}[htbp]
  \centering
  \caption{Amplitudes fitted to $568\,272$ subvolume spectra over
  $k = 0.011$--$0.135\ h\,\mathrm{Mpc}^{-1}$.}
  \label{tab:amplitude}
  \begin{tabular}{@{}llrl@{}}
    \toprule
    model fitted & subvolumes & $A$ & \\
    \midrule
    windowed   & all                & $0.9995 \pm 0.0765$ & unbiased \\
    raw theory & all                & $0.892$             & absorbs the deficit \\
    windowed   & kept by the search & $1.065$             & $0.85\sigma$ high \\
    \bottomrule
  \end{tabular}
\end{table}

\begin{figure}[htbp]
  \centering
  \includegraphics[width=\textwidth]{amplitude_fit.pdf}
  \caption{Fitting an amplitude is unbiased; searching for a realisation is not.
  With the window in the model (left) the amplitude fitted to each subvolume
  centres on $A = 1$. With the raw theory as the model (right) it centres on
  $0.892$, because that model overpredicts the data and the fit absorbs the
  difference. The outlined histogram is the one per cent of subvolumes a search
  keeps: they are the realisations that make the wrong model look right, and
  under the correct model they carry $6.5\%$ too much large-scale power.}
  \label{fig:amplitude}
\end{figure}

Put the window in the model and no search is needed: the fit returns $A = 1$.
The search instead keeps the realisations that make a misspecified model return
$A = 1$, and those realisations have $A = 1.065$ when measured correctly. An
amplitude of $1.065 \pm 0.077$ would not be reported as agreement with a
cosmology; the search manufactures precisely that and treats it as a match.

\section{What this does to every other property}
\label{sec:conditioning}

\subsection{The quantities measured}
\label{sec:quantities}

Each is a single number per subvolume, so that a selection can be asked what it
does to it. Angle brackets $\avg{\cdot}_V$ denote an average over the cells of
the subvolume, and $\delta_R$ is the density contrast smoothed with a Gaussian of
radius $R$. The transverse width of the subvolume is $\ell_\perp = L/8 =
87.5\Mpch$, which sets the smallest wavenumber it can resolve across itself,
$\Delta k_\perp = 2\pi/\ell_\perp = 0.072\ h\,\mathrm{Mpc}^{-1}$. The Nyquist
wavenumber of the grid is $k_{\rm Ny} = \pi N / L$.

\begin{description}
\item[The criterion.] The distance of a measured subvolume spectrum
  $\hat P(k)$ from a reference curve $P_{\rm ref}(k)$, over the band the
  selection uses:
  \begin{equation}
    \label{eq:criterion}
    D \;=\; \left[ \frac{1}{N_{\rm bins}}
      \sum_{k \,\le\, 2\Delta k_\perp}
      \ln^2 \frac{\hat P(k)}{P_{\rm ref}(k)} \right]^{1/2} .
  \end{equation}
  The proposal takes $P_{\rm ref} = P_{\rm th}$, the raw linear theory. The
  control of \S\ref{sec:controls} takes
  $P_{\rm ref} = P_{\rm th} \ast |\widetilde W|^2$ instead. A smaller $D$ is a
  closer match, and a selection keeps the subvolumes with the smallest $D$.

\item[Large-scale power.] The mean ratio of the measured spectrum to the
  convolved theory over the same band the criterion uses,
  $\avg{\hat P / (P_{\rm th} \ast |\widetilde W|^2)}_{k \le 2\Delta k_\perp}$.
  It is quoted against the convolved curve so that a value of one means an
  ordinary region, not a matching one.

\item[Small-scale power.] The same ratio over
  $2\Delta k_\perp < k \le 0.9\,k_{\rm Ny}$, that is, the scales the criterion
  does not look at.

\item[Tidal tensor.] $T_{ij} = \partial_i \partial_j \Phi$ with
  $\nabla^2 \Phi = \delta_R$, so in Fourier space
  \begin{equation}
    \label{eq:tidal}
    T_{ij}(\bk) \;=\; \frac{k_i k_j}{k^2}\, \delta_R(\bk) .
  \end{equation}
  Its trace is $\delta_R$, and its eigenvalues $\lambda_1 \ge \lambda_2 \ge
  \lambda_3$ describe how the region is being stretched or compressed along
  three orthogonal axes.

\item[Tidal shear.] The part of the tidal tensor that changes a region's shape
  rather than its volume, obtained by removing the trace,
  $s_{ij} = T_{ij} - \tfrac{1}{3}\delta_R \delta_{ij}$. The number quoted is its
  amplitude within the subvolume,
  \begin{equation}
    \label{eq:shear}
    s \;=\; \Big\langle s_{ij}s_{ij} \Big\rangle_V^{1/2}
      \;=\; \Big\langle T_{ij}T_{ij} - \tfrac{1}{3}\delta_R^2 \Big\rangle_V^{1/2}.
  \end{equation}
  For a Gaussian field $\avg{s_{ij}s_{ij}} = \tfrac{2}{3}\avg{\delta_R^2}$,
  which the measurement reproduces as $0.6667$.

\item[Mean density.] $\bar\delta_R = \avg{\delta_R}_V$, the smoothed density
  contrast averaged over the subvolume.

\item[Cosmic-web fractions.] Each cell is classified by how many eigenvalues of
  $T_{ij}$ are positive, a positive eigenvalue meaning collapse along that axis:
  three is a knot, two a filament, one a sheet, none a void. The four numbers
  quoted are the fractions of the subvolume's cells in each class, and they sum
  to one by construction. For a Gaussian field the expected fractions are
  $0.08$, $0.42$, $0.42$, $0.08$, and the measurement gives $0.0799$, $0.4198$,
  $0.4202$, $0.0801$.

\item[Ellipticity.] The shape of the tidal ellipsoid,
  $e = (\lambda_1 - \lambda_3)/2L$ with $L = (\sum_i \lambda_i^2)^{1/2}$. The
  usual definition divides by the trace $\delta_R$, which passes through zero
  for a region of average density and makes the ratio diverge; dividing by the
  norm keeps the same meaning and stays finite.

\item[Bulk flow.] The mean displacement of the subvolume, from the
  Zel'dovich relation $\boldsymbol{\Psi}(\bk) = i\,\bk\,\delta(\bk)/k^2$. The
  number quoted is $|\avg{\boldsymbol{\Psi}}_V|$, a vector average and then a
  magnitude, so it depends on the phases of the modes and not only on their
  amplitudes.

\item[Environment contrast.] The subvolume's mean density minus that of the
  eight subvolumes surrounding it, which asks whether the region sits in an
  overdense neighbourhood rather than whether it is overdense relative to the
  whole box.
\end{description}

Smoothing radii of $R = 20$ and $40\Mpch$ are used throughout. The first is
comfortably inside the subvolume, which is $87.5\Mpch$ across; the second is
half that width, so the smoothed field near the subvolume's faces is built
partly from material outside it, and the $R = 40\Mpch$ numbers describe the
region together with its surroundings rather than the region alone.

\subsection{The effect of selection on each}
\label{sec:each}

A region has many properties besides its power spectrum, and the criterion
mentions none of them. They move regardless, and by an amount fixed in advance.

Selection acts on one quantity. For jointly Gaussian variables the conditional
mean of any other is
$\avg{y \mid x} = \bar y + \rho\,(\sigma_y/\sigma_x)(x - \bar x)$, so in units
of each quantity's own scatter a selection that displaces $x$ by $\Delta$
displaces $y$ by
\begin{equation}
  \label{eq:rho}
  \Delta_y \;=\; \rho\,\Delta_x ,
\end{equation}
and by nothing else. No other property of the selection enters. Here the
selection displaces the subvolume's large-scale power by
$\Delta_x = 1.26\sigma$, so every other property should move by
$\rho \times 1.26$. Table~\ref{tab:rho} shows that it does.

\begin{table}[htbp]
  \centering
  \caption{Correlation with the selected quantity, the shift
  Eq.~\eqref{eq:rho} predicts, and the shift measured, keeping the closest one
  per cent of $2\,005\,008$ subvolumes. Shifts are in units of the
  region-to-region scatter.}
  \label{tab:rho}
  \begin{tabular}{lrrr}
    \toprule
    property & $\rho$ & $\rho\,\Delta_x$ & measured \\
    \midrule
    tidal shear, $R = 20\Mpch$      & $+0.451$ & $+0.57$ & $+0.62$ \\
    knot fraction, $R = 40\Mpch$    & $+0.291$ & $+0.37$ & $+0.35$ \\
    void fraction, $R = 40\Mpch$    & $+0.288$ & $+0.36$ & $+0.35$ \\
    tidal shear, $R = 40\Mpch$      & $+0.251$ & $+0.32$ & $+0.24$ \\
    filament fraction, $R = 40\Mpch$& $-0.171$ & $-0.22$ & $-0.21$ \\
    bulk flow of the region         & $+0.024$ & $+0.03$ & $+0.05$ \\
    mean density of the region      & $+0.001$ & $+0.00$ & $-0.00$ \\
    \bottomrule
  \end{tabular}
\end{table}

The last two rows carry as much weight as the first. The bulk flow and the mean
density barely correlate with the selected quantity, and they do not move.
Equation~\eqref{eq:rho} predicts the nulls as well as the effects, which is what
separates an explanation from a list of findings.

Figure~\ref{fig:summary} shows the same result graphically, with the predicted
shifts marked alongside the measured ones.

\begin{figure}[htbp]
  \centering
  \includegraphics[width=\textwidth]{selection_summary.pdf}
  \caption{Left: the convolved theory and the mean of all $2\,005\,008$
  subvolume measurements lie together well below the linear theory at large
  scales, while the selected one per cent instead follow the theory curve.
  Right: the shift in each property for the proposed criterion, for the same
  selection made against the convolved theory, and for a randomly permuted
  criterion. The short vertical marks are not fits but the shifts predicted by
  Eq.~\eqref{eq:rho}. The dotted lines mark one standard deviation of the
  natural region-to-region scatter.}
  \label{fig:summary}
\end{figure}

\section{The distortion is bounded}
\label{sec:bounded}

An obvious worry is that searching harder makes matters worse without limit. It
does not, and the reason follows from the same identity.

The displacement in the selected quantity cannot exceed what the realisation
scatter makes available. Once a search finds realisations whose subvolume power
reaches the raw theory, searching further finds realisations that match slightly
better, not realisations with more power. The displacement saturates near
$1.3\sigma$, and by Eq.~\eqref{eq:rho} everything else saturates at
$\rho \times 1.3$.

\begin{table}[htbp]
  \centering
  \caption{Shift against the number of seeds searched, keeping the best one.
  Each entry repeats the search $83\,542/N$ times; entries beyond $N = 1000$
  have too few repeats to interpret and are omitted.}
  \label{tab:saturation}
  \begin{tabular}{lrrrrr}
    \toprule
    seeds searched $N$ & 3 & 30 & 100 & 300 & 1000 \\
    \midrule
    large-scale power        & $+0.63$ & $+1.16$ & $+1.23$ & $+1.28$ & $+1.31$ \\
    tidal shear, $R = 20\Mpch$ & $+0.26$ & $+0.56$ & $+0.61$ & $+0.64$ & $+0.64$ \\
    \bottomrule
  \end{tabular}
\end{table}

A search over a thousand seeds is no worse than one over three hundred, and one
over a million would be no worse either. Figure~\ref{fig:search} shows this, and
in its right-hand panel shows something stronger: the ratio of each property's
shift to the shift in the selected quantity is flat in $N$ and equal to $\rho$.
The size of the search sets only the common factor $\Delta_x$; the structure of
the result is a property of the field.

\begin{figure}[htbp]
  \centering
  \includegraphics[width=\textwidth]{search_size.pdf}
  \caption{Left: the shift in each property against the number of seeds
  searched. Both the selected quantity and everything else saturate within a few
  hundred seeds. The shaded region has too few independent repeats to interpret.
  Right: the same shifts divided by the shift in the selected quantity, with
  each property's correlation with that quantity marked as a dashed line. The
  ratios are flat, so the size of the search changes only a common factor.}
  \label{fig:search}
\end{figure}

The same invariance holds against the tolerance rather than the search size.
Tightening the retained fraction from a quarter to two parts in ten thousand, a
range of $1250$, changes the shift in the large-scale power by $42\%$, from
$0.945\sigma$ to $1.338\sigma$, while the ratio for the tidal shear stays
between $0.43$ and $0.50$ against $\rho = 0.451$, and the ratio for the void
fraction between $0.28$ and $0.33$ against $\rho = 0.288$. Nothing in the
conclusion depends on where the threshold is placed.

Stated without inflation: the criterion biases the ensemble by a fraction of its
natural spread. A single selected region is not grotesque, and usually overlaps
the ordinary population. What is not defensible is calling such a region typical,
or treating the match as evidence that it is.

\section{Two controls}
\label{sec:controls}

Both competing explanations can be excluded by measurement rather than argument.

\emph{Is the machinery producing shifts on its own?} Replace the criterion with
a random permutation of itself, keeping everything else identical. Every
property returns $0.000 \pm 0.045$. The geometry, the estimator and the
selection arithmetic invent nothing.

\emph{Is it the act of selecting, or the curve being matched?} Change only the
target: select against the theory convolved with the window, which is the curve
the estimator is unbiased for. That target requires nothing atypical, so it
should do nothing.

\begin{table}[htbp]
  \centering
  \caption{The same selection against two different target curves, and against a
  random criterion. Keeping the closest one per cent.}
  \label{tab:controls}
  \begin{tabular}{lrrr}
    \toprule
    property & raw theory & convolved theory & random \\
    \midrule
    large-scale power             & $+1.259$ & $+0.016$ & $+0.000$ \\
    tidal shear, $R = 20\Mpch$    & $+0.617$ & $+0.056$ & $+0.000$ \\
    knot fraction, $R = 40\Mpch$  & $+0.353$ & $+0.019$ & $-0.001$ \\
    filament fraction, $R = 40\Mpch$ & $-0.205$ & $-0.013$ & $-0.001$ \\
    \bottomrule
  \end{tabular}
\end{table}

The whole effect is attributable to the choice of target curve. Selecting is not
the problem. Selecting against a curve the measurement was never going to reach
is. Figure~\ref{fig:hist} shows the same comparison as distributions, which also
distinguishes a displaced sample from a merely narrowed one.

\begin{figure}[htbp]
  \centering
  \includegraphics[width=\textwidth]{selection_histograms.pdf}
  \caption{Each property over all subvolumes and over those a selection keeps. A
  shift of the mean does not say whether the retained sample is displaced or
  merely narrowed, and the two are different for a simulation drawn from it. The
  large-scale power becomes a narrow distribution barely overlapping its parent;
  the shear and the web fractions keep most of the parent range with the centre
  moved.}
  \label{fig:hist}
\end{figure}

\section{How many independent statements these are}
\label{sec:independence}

Tables~\ref{tab:rho} and~\ref{tab:controls} list several properties, which
invites reading them as several independent confirmations. They are fewer than
they look, and the correlation matrix says how many.

The four cosmic-web fractions sum to unity at every cell, so only three of them
can be independent; measured, they sum to $1.0000$ with zero scatter, and the
correlation matrix has two exactly vanishing eigenvalues, one per smoothing
radius. The two smoothing radii largely repeat each other, the shear at
$R = 20$ and $R = 40\Mpch$ correlating at $0.76$. Of fifteen measured
quantities, three eigenvalues exceed unity and the participation ratio is
$3.49$.

The three surviving directions are interpretable, and this is the useful part.
The tidal shear correlates with the web fractions at between $-0.01$ and
$-0.09$, so the shear and the web composition are nearly orthogonal to one
another while both correlate with the selected power. They are therefore two
genuinely independent confirmations, which is stronger evidence than a dozen
correlated ones. The bulk flow is the third direction and does not move.

\begin{figure}[htbp]
  \centering
  \includegraphics[width=\textwidth]{correlations.pdf}
  \caption{Left: correlations between all measured properties. Most measure
  overlapping things, and the fifteen quantities span about three independent
  directions. Right: each property's correlation with the large-scale power the
  selection acts on, which by Eq.~\eqref{eq:rho} predicts the entire result.}
  \label{fig:corr}
\end{figure}

\section{The same argument applies to any selection rule}
\label{sec:general}

Equation~\eqref{eq:rho} says nothing about power spectra. It governs any rule
for choosing which region to resimulate.

A region chosen for a quiet merger history, for a particular halo mass, for an
isolated environment, or for resembling an observed system, is a region on which
a condition has been imposed. Every property correlated with that condition is
displaced, by the same arithmetic, whether or not the displacement was intended
and whether or not it is reported. The choice of region is part of the numerical
experiment and belongs in its interpretation, in the same way a selection
function belongs in the interpretation of a survey.

The useful half is that this is predictable before anything is run. The
correlations are properties of the field and the geometry, not of the selection,
and a few hundred realisations measure them. Anyone proposing a criterion can be
told in advance which properties of the resimulated region it will distort and
by how much.

The converse is equally useful. A property uncorrelated with the criterion is
untouched: the mean density of the region has $\rho = 0.001$ with the subvolume's
large-scale power and shows no shift at any search size or tolerance. A criterion
is safe for exactly those quantities it does not correlate with, and the
correlation matrix identifies them.

\section{Recommendations}
\label{sec:recommendations}

\begin{enumerate}
\item \textbf{Compare in the basis you measured in.} Convolve the theory with
  the subvolume's window and compare against that. It reproduces the measurement
  to a fraction of a per cent, and requires no search: a typical region already
  agrees with it.

\item \textbf{If a number is wanted, fit one.} Put the window in the model and
  fit an amplitude. It returns $A = 1.000 \pm 0.077$ with no seed selection at
  all, and the width of that interval is the honest statement of how well a
  region of this size constrains the amplitude.

\item \textbf{Do not treat agreement with the unwindowed theory as validation.}
  A region that agrees with it is, for that reason, not a representative region.

\item \textbf{If a region must be chosen on some criterion}, and often one must
  be, measure the correlation between that criterion and the properties the
  science depends on, and quote the induced shift with the result. That converts
  a hidden bias into a stated systematic.
\end{enumerate}

\section{Method}
\label{sec:method}

Realisations are generated with a monofonIC fork that writes the Lagrangian
density field directly, avoiding the particle stage. Each realisation is a
$128^3$ grid in a $700\Mpch$ box at $z = 200$ with second-order Lagrangian
perturbation theory and no amplitude fixing. Twenty-four subvolumes are measured
per realisation, drawn from the $192$ that tile the box in three orientations.

The power spectrum of a subvolume is measured by masked fast Fourier transform
with the normalisation of \citet{Park1994}, and the theory is convolved with the
same window on the same grid, so measurement and prediction share every
convention.

The tidal tensor is $T_{ij} = \partial_i \partial_j \Phi$ with
$\nabla^2 \Phi = \delta$, evaluated on the Gaussian-smoothed field. The shear is
its traceless part, and for a Gaussian field
$\avg{s^2} = \tfrac{2}{3}\avg{\delta^2}$, which the measurement reproduces as
$0.6667$. Cells are classified by the number of positive eigenvalues of $T$; for
a Gaussian field the expected fractions are $0.08$, $0.42$, $0.42$ and $0.08$
for knots, filaments, sheets and voids, and the measurement gives $0.0799$,
$0.4198$, $0.4202$ and $0.0801$.

Uncertainties come from a bootstrap over whole realisations, never over
subvolumes, because subvolumes within one realisation share that realisation's
modes. At the one per cent tolerance the retained subvolumes come from
$17\,835$ distinct realisations, $1.12$ per realisation, so they are close to
independent draws in any case.

{\sloppy
Code, run provenance and the analysis are in the \texttt{cosmo-sim} repository
under \texttt{scripts/ic\_search/}; the seed ranges and configuration of each
production run are recorded in \texttt{scripts/ic\_search/RUNS.md}.
\par}

\begin{thebibliography}{9}
\bibitem[Park et al.(1994)]{Park1994} Park, C., Vogeley, M.~S., Geller, M.~J.,
  \& Huchra, J.~P. 1994, ApJ, 431, 569
\end{thebibliography}

\end{document}
